% Generated draft: requires mathematical and source review.
\documentclass[11pt]{article}
\usepackage[margin=25mm]{geometry}
\usepackage{fontspec}
\setmainfont{Latin Modern Roman}
\usepackage{amsmath,amssymb,mathtools,unicode-math}
\setmathfont{Latin Modern Math}
\usepackage{enumitem,needspace,xurl,hyperref}
\setlength{\parindent}{0pt}
\setlength{\parskip}{5pt}
\setlength{\emergencystretch}{3em}
\newcommand{\N}{\mathbb{N}}
\newcommand{\Z}{\mathbb{Z}}
\newcommand{\Q}{\mathbb{Q}}
\newcommand{\R}{\mathbb{R}}
\newcommand{\C}{\mathbb{C}}
\newcommand{\F}{\mathbb{F}}
\newcommand{\A}{\mathbb{A}}
\newcommand{\PP}{\mathbb{P}}
\newcommand{\E}{\mathbb{E}}
\newcommand{\eps}{\varepsilon}
\newcommand{\dd}{\,\mathrm{d}}
\begin{document}
\section*{UnsolvedMath ID: 801024}
\textbf{Source problem code:} catalogue entry
\subsection{Definitions and mathematical statement}
A profinite group is a topological group isomorphic, as a topological group, to an inverse limit of finite discrete groups. Let \(\mathcal P\) denote the class of profinite groups, and let \(L_{\mathrm{grp}}=\{\cdot,{}^{-1},1\}\) be the first-order language of groups. For an \(L_{\mathrm{grp}}\)-sentence \(\varphi\), define the class
\[
\mathcal C_\varphi=\{G\in\mathcal P:G\models\varphi\},
\]
where satisfaction is evaluated on the underlying group; the language contains no predicate for the topology.

The supplied source asks: what complexities can \(\mathcal C_\varphi\) have as \(\varphi\) ranges over first-order group sentences? The restriction to profinite groups is indicated by the surrounding source discussion, although it is omitted from the displayed class in the source question. No restriction to metrizable, countably based, or finitely generated profinite groups is supplied.

This is an unresolved formulation of a question, rather than a specified complexity assertion: the source does not define a notion of complexity, a parameter space or encoding of profinite groups, or an associated topology, measurable structure, or effective presentation. Thus \(\mathcal C_\varphi\) is presently a class, and no particular descriptive or computational complexity classification can be stated without additional choices.

As context, the supplied text attributes to Nies, Segal, and Tent (2021), reference label \([\mathrm{NST}]\), a study of first-order expressiveness for profinite groups, and reports that many such groups admit axiomatisation by a single sentence relative to this reference class. This contextual report is retained as supplied; it does not specify the intended complexity framework.
\subsection{Short English statement}
For a first-order sentence in the language of groups, consider the profinite groups satisfying it. The question asks how complex this class can be. A precise classification problem remains unspecified because the source gives neither a complexity notion nor a way to encode the groups.
\subsection{Sources}
\begin{itemize}
\item Local UnsolvedMath registry, permanent ID \(801024\), source problem code \texttt{catalogue entry}, titled \emph{Open Problems in Mathematical Logic --- Problem 36}. The supplied registry labels the problem \texttt{open}; this is a supplied status label, not an independently verified mathematical assessment.
\item Barmpalias, Gao et al., \emph{Open Problems in Mathematical Logic} (2026), Problem 36, \url{https://arxiv.org/abs/2608.26628}. The registry attributes the entry to Barmpalias, Gao et al. and gives proposed year 2026. The supplied provenance records access on 2026-10-16 and extraction by \texttt{arxiv\_latex}; these are supplied provenance statements, not a claim of independent source access or verification. The supplied upstream revision is \texttt{d371ee917ac9647fbb06cd6fc529deec6f3dda82}.
\item Nies, Segal, and Tent (2021), reference label \([\mathrm{NST}]\), cited in the supplied problem text for work on first-order expressiveness for profinite groups. No title, publication details, or link for this reference is supplied.
\end{itemize}
\end{document}
